%%%%%%%%%%%%%%%%%%%%%%%%%%%%%%%%%%%%%%%%%%%%%%%%
% COPYRIGHT: (C) 2012-2026 FAU FabLab and others
% Lizenz: CC-BY-SA 3.0 (wie alle anderen Einweisungen des FAU FabLab)
%%%%%%%%%%%%%%%%%%%%%%%%%%%%%%%%%%%%%%%%%%%%%%%%


\PassOptionsToPackage{english}{babel} % zusätzlich zu ngerman aus fablab-document
\newcommand{\basedir}{fablab-document}
\documentclass[13pt]{\basedir/fablab-document}
\usepackage{fancybox} %ovale Boxen für Knöpfe
\usepackage{amssymb} % Symbole für Knöpfe
\usepackage{subfigure,caption}
\usepackage{eurosym}
\usepackage{tabularx} % Tabellen mit bestimmtem Breitenverhältnis der Spalten
\usepackage{wrapfig} % Textumlauf um Bilder
\renewcommand{\texteuro}{\euro}
\usepackage{ifthen}
\usepackage{xspace}
\def\tabularnewcol{&\xspace} % hässlicher Workaround von http://tex.stackexchange.com/questions/7590/how-to-programmatically-make-tabular-rows-using-whiledo


\usepackage{tabularx} % Tabelle mit teilweise gleich großen Spalten
\title{Instruction list general workshop rules}
\fancyfoot[C]{Einweisungsliste Werkstatt en.tex}
\fancyfoot[L]{Instruction list no. \underline{\hspace{3em}}}

\begin{document}
\selectlanguage{english}
%\maketitle

With my signature I bindingly confirm that

\begin{itemize}
\item I have read and understood the general workshop rules (code of conduct, traffic light system, electronics)
\item there was an opportunity for questions and discussion
\item I had the opportunity to obtain a copy of the rules
\end{itemize}


% einfach kopiert von Einweisungsliste Lasercutter

\newcounter{i}
\setcounter{i}{1}

\newcommand{\leerezeile}{\hspace{2em} \tabularnewcol \hspace{3em} \tabularnewcol \hspace{2.5em} \tabularnewcol \hspace{2.5em} \tabularnewcol \vbox{\vspace{2em}} \tabularnewcol \tabularnewcol \tabularnewcol \tabularnewline \hline}

\begin{tabularx}{\textwidth}{|l|l|l|l|X|X|X|X|}
  \hline
  \textbf{No.} & \textbf{Date} & \textbf{from} & \textbf{to} & \textbf{Name} & \textbf{Signature} & \textbf{Instructor} & \textbf{Signature} \\ \hline
  \whiledo{\value{i}<14}%
  {%
    \stepcounter{i} \leerezeile
  }%
  \leerezeile % doofer Workaround, eigentlich sollte das auch in der Forschleife gehen! Ohne dies wird die Spaltenbegrenzung von Spalte 1 zu weit gezeichnet.
\end{tabularx}

\end{document}